% Big Data in Omics and Imaging 中文工程共享导言区。
\usepackage{amsmath,amssymb,amsthm,mathtools,bm}
\usepackage{tikz}
\usetikzlibrary{arrows.meta,positioning,calc,backgrounds,shapes.geometric}
\usepackage{graphicx,booktabs,array,tabularx,enumitem,xcolor,caption,fancyhdr,placeins,needspace}
\usepackage[a4paper,top=2.6cm,bottom=2.6cm,left=2.6cm,right=2.6cm,headsep=0.7cm,footskip=1.0cm]{geometry}
\setmainfont{TeX Gyre Termes}
\definecolor{inkblue}{RGB}{28,62,110}
\definecolor{plotgray}{RGB}{35,31,32}
\linespread{1.32}
\setlength{\parskip}{0.25em}
\setlength{\emergencystretch}{1.2em}
\tolerance=2000
\setlength{\headheight}{15pt}
\ctexset{
 chapter={format=\centering\heiti\zihao{2},name={第,章},number=\chinese{chapter},beforeskip=1.2cm,afterskip=1.2cm},
 section={format=\heiti\zihao{3},beforeskip=1.1ex plus .3ex minus .2ex,afterskip=.8ex plus .2ex},
 subsection={format=\heiti\zihao{4},beforeskip=1.0ex plus .3ex minus .2ex,afterskip=.6ex plus .2ex},
 figurename=图,tablename=表
}
\captionsetup{font=small,labelsep=quad,labelfont=bf,skip=8pt}
\captionsetup[figure]{justification=justified,singlelinecheck=false}
\captionsetup[table]{justification=justified,singlelinecheck=false}
\pagestyle{fancy}
\fancyhf{}
\fancyhead[LE]{\small\kaishu\leftmark}
\fancyhead[RO]{\small\kaishu\rightmark}
\fancyfoot[C]{\small\thepage}
\renewcommand{\headrulewidth}{0.4pt}
\fancypagestyle{plain}{\fancyhf{}\fancyfoot[C]{\small\thepage}\renewcommand{\headrulewidth}{0pt}}
\makeatletter
\renewcommand{\chaptermark}[1]{\markboth{\CTEXthechapter\quad #1}{}}
\renewcommand{\sectionmark}[1]{\markright{\thesection\quad #1}}
\makeatother
% 通用定理环境；本书印刷编号由book*环境显式保留。
\newtheoremstyle{cnprinciple}{7pt}{7pt}{\normalfont}{} {\bfseries}{}{0.7em}{}
\theoremstyle{cnprinciple}
\newtheorem{resultshared}{}[chapter]
\newtheorem{principle}[resultshared]{原理}
\newtheorem{definition}[resultshared]{定义}
\newtheorem{theorem}[resultshared]{定理}
\newtheorem{lemma}[resultshared]{引理}
\newtheorem{corollary}[resultshared]{推论}
\newtheorem{proposition}[resultshared]{命题}
\newtheorem{remark}[resultshared]{注}
\newtheorem{example}[resultshared]{例}
\newtheorem{assumption}[resultshared]{假设}
\newtheorem{conjecture}[resultshared]{猜想}
\newtheorem{codesnippet}[resultshared]{代码片段}
\newtheorem{algorithm}[resultshared]{算法}
\newtheorem{problem}{习题}[chapter]
\renewcommand{\proofname}{证明}
\newenvironment{remarkenv}{\par\noindent\textbf{注。}\ignorespaces}{\par}
\newcommand{\PA}{\mathrm{PA}}
\newcommand{\pa}{\mathrm{pa}}
\newcommand{\ND}{\mathrm{ND}}
\newcommand{\AN}{\mathrm{AN}}
\newcommand{\DE}{\mathrm{DE}}
\newcommand{\EE}{\mathbb{E}}
\newcommand{\RR}{\mathbb{R}}
\DeclareMathOperator{\Var}{Var}
\DeclareMathOperator{\Cov}{Cov}
\DeclareMathOperator{\tr}{tr}
\DeclareMathOperator{\diag}{diag}
\DeclareMathOperator{\supp}{supp}
\DeclareMathOperator{\sign}{sign}
\usepackage{listings}
\lstset{basicstyle=\ttfamily\footnotesize,columns=fullflexible,keepspaces=true,
 breaklines=true,numbers=left,numberstyle=\tiny\color{gray},numbersep=6pt,
 frame=single,framerule=.3pt,rulecolor=\color{gray!50},showstringspaces=false,
 tabsize=2,aboveskip=6pt,belowskip=6pt}
\usepackage[most]{tcolorbox}
\newtcolorbox{sourcebox}[1][]{enhanced,breakable,colback=gray!3,colframe=gray!50,
 boxrule=.4pt,arc=0pt,left=7pt,right=7pt,top=5pt,bottom=5pt,title={#1},coltitle=black,
 colbacktitle=gray!10,fonttitle=\bfseries}
\usetikzlibrary{decorations.pathreplacing,decorations.markings,fit,patterns}
\tikzset{
 vtxb/.style={vtx,very thick},
 vtxs/.style={vtx,minimum size=5.5mm,font=\footnotesize},
 lat/.style={latent},
 edb/.style={->,very thick},
 bid/.style={<->,semithick},
 bidd/.style={<->,semithick,dashed}
}

\newcommand{\indep}{\mathrel{\perp\!\!\!\perp}}
\newcommand{\nindep}{\mathrel{\not\!\perp\!\!\!\perp}}
\DeclareMathOperator*{\argmax}{arg\,max}
\DeclareMathOperator{\doop}{do}
\newcommand{\dox}[1]{\mathrm{do}(#1)}
\def\doxp(#1){\mathrm{do}(#1)}
\tikzset{
 causal/.style={>={Stealth[length=2.1mm,width=1.5mm]},semithick},
 vtx/.style={circle,draw,inner sep=1pt,minimum size=7mm,font=\small},
 latent/.style={vtx,dashed},
 ed/.style={->,semithick},
 edd/.style={->,semithick,dashed}
}
\setlist[enumerate]{leftmargin=2.3em,itemsep=0.35em,topsep=0.4em}
\renewcommand{\arraystretch}{1.3}
\setcounter{secnumdepth}{2}
\setcounter{tocdepth}{2}
\usepackage{longtable,multicol}
\usepackage{url}
\usepackage[bookmarks=true,bookmarksnumbered=true,unicode=true,colorlinks=true,linkcolor=inkblue,citecolor=inkblue,urlcolor=inkblue]{hyperref}

\usepackage[most]{tcolorbox}
\usepackage{longtable,multirow,adjustbox}
\newcommand{\Tr}{\operatorname{Tr}}
\newcommand{\soft}{\operatorname{soft}}
\newcommand{\argmin}{\operatorname*{arg\,min}}
\newcommand{\sgn}{\operatorname{sgn}}
\newcommand{\vecop}{\operatorname{vec}}
\newcommand{\E}{\mathbb{E}}
\newcommand{\R}{\mathbb{R}}
\newcommand{\origfig}[2]{\setcounter{figure}{\numexpr#1-1\relax}\caption{#2}}
\newcommand{\origtable}[2]{\setcounter{table}{\numexpr#1-1\relax}\caption{#2}}
\newenvironment{bookexample}[1]{\par\medskip\noindent\textbf{例 #1}\quad\ignorespaces}{\par\medskip}
\newenvironment{booktheorem}[1]{\par\medskip\noindent\textbf{定理 #1}\quad\ignorespaces}{\par\medskip}
\newenvironment{bookdefinition}[1]{\par\medskip\noindent\textbf{定义 #1}\quad\ignorespaces}{\par\medskip}
\newenvironment{booklemma}[1]{\par\medskip\noindent\textbf{引理 #1}\quad\ignorespaces}{\par\medskip}
\newenvironment{bookalgorithm}[1]{\par\medskip\Needspace{5\baselineskip}\noindent\textbf{算法 #1}\quad\ignorespaces}{\par\medskip}
\newenvironment{bookproof}{\par\noindent\textbf{证明}\quad\ignorespaces}{\hfill$\square$\par\medskip}
\tikzset{vtxb/.style={vtx,very thick},lat/.style={vtx,dashed},edb/.style={->,very thick},bid/.style={<->,dashed,semithick},edd/.style={->,dashed,semithick}}
\setcounter{secnumdepth}{3}
\setcounter{tocdepth}{3}
\hypersetup{pdftitle={组学与影像大数据：整合分析与因果推断（第1章中文样章）},pdfauthor={Momiao Xiong}}

% 本书第五、六级印刷标题；原号显式给出，目录与书签仍可定位。
\newcommand{\bookdeepheading}[3]{\par\Needspace{4\baselineskip}\phantomsection
 \addcontentsline{toc}{#1}{\protect\numberline{#2}#3}
 \medskip\noindent\textbf{#2\quad #3}\par\nobreak\smallskip\nobreak}
\setcounter{tocdepth}{6}
\makeatletter
\renewcommand*\l@paragraph{\@dottedtocline{4}{10em}{6.2em}}
\renewcommand*\l@subparagraph{\@dottedtocline{5}{13em}{7em}}
\makeatother

% Explicit seventh-level book headings.
\makeatletter
\providecommand*\l@booklevelseven{\@dottedtocline{6}{16em}{8em}}
\providecommand*\toclevel@booklevelseven{6}
\makeatother
\hypersetup{bookmarksdepth=6}
